\section{2019 转向：从 AI 音箱到 AR 眼镜}

上一章讲 Rokid 的入口判断来源，本章进入最关键的战略转向。Rokid 从第一天起就认为 AI 和 AR 会变成一件事，但早期 AR 硬件太重、太怪、无法产品化，于是先用 AI 消费硬件和智能音箱训练语音、供应链、生产、销售和用户反馈能力。2019 年，当大厂以补贴和平台逻辑进入智能音箱，Rokid 判断音箱不可能成为真正入口，于是一周内从 AI 赛道切换到 AR 赛道。

这个转向很痛苦，但不是突然改方向。祝铭明强调，内部一直是 AI 和 AR 两条腿并行，只是原来 AI 产品线可以以战养战；当音箱竞争变成不健康烧钱，继续投入就不再合理。Rokid 要做的是未来战争，而不是在一个不属于自己的平台战场上消耗现金。

\begin{figure}[H]
\centering
\includegraphics[width=0.94\textwidth]{ar-pivot-decision.png}
\caption{2019：从 AI 音箱转向 AR，一周内切换赛道，代价是组织重构。自制概念图，依据 00:31:41--00:59:17 对谈内容整理。}
\end{figure}

\begin{knowledgebox}{读图：不是 AI 放弃，而是载体切换}
左侧 AI 音箱是相对成熟的语音产品，能训练团队，但输入输出窄、场景不碎片、平台属性弱。右侧 AR 眼镜更难，但更接近随身 AI 入口。Rokid 的转向是从短期可做的产品，切到长期更大的入口。
\end{knowledgebox}

\subsection{两条腿并行：AI 产品线与 AR 实验室}

本节补上一个容易被误读的细节。祝铭明说，Rokid 不是突然从 AI 改做 AR，而是从一开始就把 AI 和 AR 看成会合流的一件事。早期 AR 太重、太像实验装置，没法成为大众产品；但 NLP、ASR、TTS 等语音技术已经可以产品化，于是 AI 硬件先承担两个任务：一是让公司和市场保持接触，二是训练硬件供应链、生产、销售和用户反馈能力。

这也是“以战养战”的真正含义。音箱不是最终平台，但它能让一个软件/系统出身的团队学会做硬件。祝铭明提到，2014 到 2016 年其实是在交硬件学费：产品定义、定价、供应链、消费市场接受度都要被市场检验。Rokid 早期的 Alien 陪伴机器人很酷，但过于超前；最后被简化成音箱，门槛也随之降低，巨头补贴进入后，创业公司继续烧钱就不再合理。

\noindent\begin{tabularx}{\textwidth}{p{0.18\textwidth}p{0.32\textwidth}X}
\toprule
路线 & 当时的作用 & 后来的问题或价值 \\
\midrule
Alien 陪伴机器人 & 把摄像头、表情、语音和陪伴放进一个产品 & 概念超前，但价格和消费市场接受度不足。 \\
AI 音箱 & 训练语音产品、生产、销售和供应链 & 门槛降低后进入补贴竞争，平台属性不足。 \\
AR Lab & 长期预研光波导、微显示、成像等技术 & 让 2019 转向不是从零开始。 \\
智能眼镜 & 重新承载随身 AI 入口假设 & 更难，但平台上限更高。 \\
\bottomrule
\end{tabularx}

\begin{knowledgebox}{课堂提示：先做产品不等于放弃平台理想}
Rokid 先做音箱，是因为语音 AI 更早能落地；它后来放弃音箱，是因为音箱无法支撑平台假设。这个顺序说明，创业公司的阶段性产品可以服务长期能力建设，但不能反过来绑架长期方向。
\end{knowledgebox}

\subsection{音箱为什么不是入口}

本节把音箱的失败逻辑讲清楚。祝铭明认为，入口产品需要足够宽的输入输出通道、足够多的使用场景和足够高的碎片化频次。音箱的通道窄，主要是读书、音乐、简单问答；场景也被限制在客厅和卧室，用户醒着的大部分时间并不在这些固定场景。它可以是一个不错的产品，但很难成为手机级入口。

大公司把智能音箱当平台做，补贴、烧钱、争生态，这对创业公司不利。Rokid 不把音箱视为平台，所以果断退出。这个判断很重要：创业公司不能只看“市场热不热”，还要判断这个热度是不是符合自己的平台假设。

\begin{warningbox}{平台幻觉}
一个产品有 AI、有语音、有硬件，不代表它是平台。平台必须有高频入口、足够宽的信息通道、可持续生态和用户默认使用习惯。缺少这些条件，烧钱也可能只是在放大一个普通产品。
\end{warningbox}

\begin{knowledgebox}{入口判断清单}
\noindent\begin{tabularx}{\textwidth}{p{0.22\textwidth}p{0.32\textwidth}X}
\toprule
判断项 & 音箱的问题 & 眼镜的可能性 \\
\midrule
输入输出宽度 & 主要依赖语音，输出也偏串行 & 语音、视觉、显示和环境感知可以叠加。 \\
碎片化频次 & 固定在卧室、客厅等有限场景 & 戴在身上，可覆盖路上、店里、会议、旅行。 \\
替代手机路径 & 许多任务手机已经足够方便 & 看时间、扫码、翻译、通知等短任务可缩短路径。 \\
生态空间 & 内容类型集中在音乐、读书和问答 & 若使用时长和显示成立，应用生态更宽。 \\
创业公司优势 & 巨头可以补贴和铺量 & 前沿工艺和产品定义尚未标准化。 \\
\bottomrule
\end{tabularx}
\end{knowledgebox}

\subsection{裁员与组织重构}

本节从产品判断进入组织代价。战略转向的代价是组织重构。Rokid 当时有两栋楼，一栋做音箱，一栋做 AR；转型后音箱楼基本清空，裁掉超过一半的人。祝铭明强调痛苦不在方向判断本身，而在于许多优秀员工因为方向变化被迫离开。公司当时给出 N+1，并留下未来 AR 跑顺后优先回来的口子。

这段对硬件创业很有启发。软件产品转向可以更轻，硬件产品线转向往往牵动供应链、团队技能、库存、生产计划和现金流。Rokid 能转过去，是因为账上还有钱、投资人理解长期方向、AR Lab 已经提前跑了几年。如果等现金彻底耗尽，转型窗口就会消失。

\begin{importantbox}{转型窗口}
硬件公司做大转向，必须在现金还足够、团队还有能力、下一条路线已有预研时启动。等产品线彻底失败再转，组织、供应链和投资人信心通常已经不足。
\end{importantbox}

\subsection{疫情中的第一次 PMF}

接下来要看的是转型后的偶然性和必然性。Rokid 2019 年 10 月做调整，2020 年 1 月中旬召开眼镜产品招商会；几天后疫情爆发，线下流动被限制。由于眼镜可以配合红外传感做非接触测温，第一代较重的 B 端产品反而迅速产业化，卖到多个国家。祝铭明没有把这称为“感谢疫情”，但承认它客观上加速了产品、供应链和收入曲线。

这段最有价值的地方，是区分 PMF 与运气。疫情是外部事件，Rokid 无法预测；但如果 2019 年没有提前切到 AR、没有提前开招商会、没有供应链准备，就无法接住这个事件。硬件创业的 PMF 往往不是一个纯粹的产品灵感，而是路线、时点、供应链和场景突然对齐。

\begin{knowledgebox}{实践经验：偶然性需要被准备接住}
祝铭明说，如果招商会再晚五天，结果可能完全不同。这不是鼓励迷信时间点，而是提醒：前沿硬件公司要在窗口打开前完成足够多的准备。机会出现时，能不能交付往往比能不能讲故事更关键。
\end{knowledgebox}

\subsection{本章小结}

Rokid 的 2019 转向不是从 AI 到非 AI，而是从“语音 AI 产品”转向“AI 的随身入口”。这次转向训练了公司对平台属性、供应链和组织代价的判断，也说明硬件转型要在资金、预研、组织和外部时点之间找到窗口。

